%% kapa: exam=1 round=29 year=2018 period=1 form=A subject=real_estate subject_ko=부동산학원론 seq=5 number=85
%% source: Q-Net 감정평가사 기출문제 (cst00309) artlSeq=5202808 file=extracted/1008/5202808/2018년 제29회 감정평가사 1차시험 기출문제/2018년 제29회 감정평가사 1차시험 기출문제/1교시 A형.pdf
%% conversion: kordoc 4.24.51 -> tools/gen1.py (auto); review=proofread-sonnet; flags=none
\begin{question}{85}
부동산과 준부동산에 관한 설명으로 옳은 것은? (다툼이 있으면 판례에 따름)
\begin{choices}
  \item 신축 중인 건물은 사용승인이 완료되기 전에는 토지와 별개의 부동산으로 취급되지 않는다.
  \item 개개의 수목은 명인방법을 갖추더라도 토지와 별개의 부동산으로 취급되지 않는다.
  \item 토지에 정착된 담장은 토지와 별개의 부동산으로 취급된다.
  \item 자동차에 관한 압류등록은 자동차 등록원부에 한다.
  \item 총톤수 10톤 이상의 기선(機船)과 범선(帆船)은 등기가 가능하다.
\end{choices}
\end{question}
\answer{4}
